\documentclass[preprint,12pt]{elsarticle}

\usepackage{amsmath,amssymb}
\usepackage{graphicx}
\usepackage{siunitx}
\usepackage{booktabs}
\usepackage{hyperref}

\begin{document}

\begin{frontmatter}

\title{Shape optimisation of acoustic beacons in bat-pollinated flowers: a boundary element method study}

\author[1]{Ant Newman}
\affiliation[1]{organization={TortoiseAI / University of Bristol (affiliation TBD)}}

\begin{abstract}
Several bat-pollinated plant species have evolved concave structures that function as acoustic retroreflectors, preferentially reflecting echolocation energy back toward approaching pollinators. Empirical characterisation of these reflectors is mature, but no formal scattering model or shape optimisation has been published. We present the first open-source boundary element method (BEM) solver applied to floral acoustic reflectors, and the first formal shape optimisation of structures evolved to attract echolocating bats. The Helmholtz problem is solved for rigid sound-hard scatterers using the Burton-Miller combined field integral equation. Parametric reflector families (spherical cap, ellipsoidal cap, truncated Chebyshev profile) are optimised against biologically meaningful objective functions using CMA-ES (single-objective) and NSGA-II (multi-objective). We compare the Pareto frontier of acoustic conspicuousness versus surface area against the geometries of \textit{Marcgravia evenia} and \textit{Mucuna holtonii}. [Results paragraph to be added once production optimisation runs complete.]
\end{abstract}

\begin{keyword}
acoustic scattering \sep boundary element method \sep echolocation \sep shape optimisation \sep biomimetics \sep bat pollination
\end{keyword}

\end{frontmatter}

\section{Introduction}
\label{sec:introduction}

Acoustic coevolution between bat-pollinated flowers and their nectar-feeding pollinators has received substantial empirical attention. Two canonical systems anchor the literature: the concave vexillum of \textit{Mucuna holtonii}, which functions as a cat's-eye retroreflector across a wide range of incidence angles~\cite{vonHelversen1999,vonHelversen2003jcp}, and the concave dish-leaf of \textit{Marcgravia evenia}, which halves the foraging time required for its glossophagine pollinator to locate flowers~\cite{Simon2011}. A third example, the absorptive cephalium of \textit{Espostoa frutescens}, produces a frequency-specific attenuation notch tuned to the echolocation band of its pollinator~\cite{Simon2023}.

Mathematical modelling has lagged the empirical work. Only one published numerical scattering computation for any floral reflector exists: the finite element (FEM) study of Simon et al.~\cite{Simon2020}, which characterised spherical caps at three radii and three depths. No boundary element computation has been applied to this class of scatterer, and no systematic shape optimisation has been performed against biologically meaningful objectives. The Helmholtz number regime relevant to these reflectors (He $\approx 3$ to $30$) places them in the resonance-to-optical transition, where full-wave methods are required and ray-tracing approximations are inadequate.

This paper addresses the gap by presenting a validated open-source BEM solver (\texttt{acoustic-beacon-optimiser}) applied to parameterised concave reflectors, coupled to derivative-free optimisers for both single-objective and multi-objective settings. The core scientific question is whether the natural reflector geometries are Pareto-optimal in the trade-off between acoustic conspicuousness and structural cost, or whether they deviate from the frontier in ways that reveal additional constraints not captured in the purely acoustic model.

\section{Methods}
\label{sec:methods}

\subsection{Forward problem: Helmholtz BEM}
\label{sec:bem}

We solve the exterior Helmholtz problem for a rigid (sound-hard) scatterer. Given an incident plane wave $p_{\mathrm{inc}}(\mathbf{x}) = \exp(i k \hat{\mathbf{n}} \cdot \mathbf{x})$ with wavenumber $k = 2\pi f / c$ and direction $\hat{\mathbf{n}}$, the total field $p = p_{\mathrm{inc}} + p_{\mathrm{scat}}$ satisfies
\begin{align}
\nabla^2 p + k^2 p &= 0 \quad \text{in the exterior domain}, \\
\partial p / \partial n &= 0 \quad \text{on the reflector surface}, \\
\end{align}
together with the Sommerfeld radiation condition. To avoid spurious interior resonances we use the Burton-Miller combined field integral equation:
\begin{equation}
\left( \tfrac{1}{2} I + D_k + \alpha H_k \right) \phi = - \left( p_{\mathrm{inc}} + \alpha \, \partial p_{\mathrm{inc}} / \partial n \right),
\end{equation}
with coupling parameter $\alpha = i/k$, where $D_k$ is the double-layer operator and $H_k$ is the hypersingular operator. The implementation uses \texttt{bempp-cl} 0.4~\cite{bempp}. Surface meshes are generated by Gmsh~\cite{gmsh} via Boolean operations (sphere cut by half-space for caps, direct structured triangulation for Chebyshev profiles).

\subsection{Reflector parameterisations}

Three axisymmetric families are defined:
\begin{description}
\item[Spherical cap] (2 parameters) with radius $r$ and half-angle $\theta$; depth $d = r(1 - \cos\theta)$, aperture $a = r \sin\theta$.
\item[Ellipsoidal cap] (3 parameters) with radial semi-axis $a$, axial semi-axis $c$, and half-angle $\theta$; reduces to a spherical cap when $a = c$.
\item[Chebyshev profile] (N parameters) with axial profile $z(\rho) = \sum_{n=0}^{N-1} c_n T_n(2\rho/a - 1)$ evaluated over $\rho \in [0, a]$.
\end{description}

\subsection{Objective functions}

The monostatic target strength is
\begin{equation}
\mathrm{TS}(f, \theta) = 10 \log_{10} \left( |p_{\mathrm{scat}}|^2 r^2 / |p_{\mathrm{inc}}|^2 \right),
\end{equation}
evaluated in the far field in the backscatter direction. Integrated conspicuousness is defined as the call-spectrum-weighted and angle-integrated backscatter intensity:
\begin{equation}
\mathrm{IC} = 10 \log_{10} \left[ \int_{f_{\min}}^{f_{\max}} W(f) \int_0^{\theta_{\max}} 10^{\mathrm{TS}(f,\theta)/10} \sin\theta \, d\theta \, df \right],
\end{equation}
where $W(f)$ is the normalised power spectrum of the bat call. We use published spectra for \textit{Glossophaga soricina}~\cite{Simon2006} and \textit{Leptonycteris yerbabuenae}~\cite{GonzalezTerrazas2016}. Surface area $\mathrm{SA}$ is computed from the triangle-mesh connectivity.

\subsection{Optimisation}

Single-objective optimisation uses CMA-ES~\cite{Hansen2001} via the \texttt{cma} package in an ask-and-tell driver, maximising IC subject to a surface-area penalty. Multi-objective optimisation uses NSGA-II~\cite{Deb2002} via \texttt{pymoo}~\cite{pymoo}, producing the Pareto frontier of $(\mathrm{IC}, \mathrm{SA})$.

\subsection{Validation}

The BEM solver is validated against the analytical Mie series for a rigid sphere at Helmholtz numbers $\mathrm{He} \in \{5, 10, 20\}$. Results are reported in Figure~\ref{fig:mie}. Cross-validation against the FEM results of Simon et al.~\cite{Simon2020} is deferred pending access to the published numerical data.

\section{Results}
\label{sec:results}

\subsection{Mie-series benchmark}

Figure~\ref{fig:mie} shows the absolute error between BEM-computed and Mie-series monostatic target strengths for a rigid sphere. Agreement is within \SI{0.03}{\decibel} at He = 5 and He = 10 using approximately 10 elements per wavelength.

\begin{figure}[h]
\centering
\includegraphics[width=0.7\linewidth]{figures/fig01_mie_validation.png}
\caption{BEM-versus-Mie backscatter TS agreement for a rigid unit sphere. Dashed line marks the \SI{0.5}{\decibel} tolerance used in the regression test suite.}
\label{fig:mie}
\end{figure}

\subsection{Marcgravia evenia approximation}

Figure~\ref{fig:marcgravia} shows the spectral directivity of a spherical-cap approximation to the \textit{Marcgravia evenia} dish-leaf.

\begin{figure}[h]
\centering
\includegraphics[width=0.7\linewidth]{figures/fig02_marcgravia_ts_heatmap.png}
\caption{Monostatic target strength of a spherical-cap approximation to \textit{Marcgravia evenia} (aperture \SI{17.5}{\milli\metre}, depth \SI{25}{\milli\metre}).}
\label{fig:marcgravia}
\end{figure}

\subsection{Single-objective optimum}

[To be completed from notebook 03 production run.]

\subsection{Pareto frontier}

[To be completed from notebook 04 production run.]

\section{Discussion}
\label{sec:discussion}

[Placeholder: interpret whether natural geometries sit on or inside the Pareto frontier. If on, acoustic optimality is plausible; if inside, identify candidate additional constraints (structural, developmental, phylogenetic) and discuss which can be tested empirically.]

\section{Conclusions}
\label{sec:conclusions}

[Placeholder.]

\section*{Code and data availability}

The computational methods described here are implemented in the open-source Python package \texttt{acoustic-beacon-optimiser}, available at \url{https://github.com/antnewman/acoustic-beacon-optimiser} under the MIT licence. All figures in this paper are reproducible from the Jupyter notebooks in the repository's \texttt{notebooks/} directory. Raw optimisation outputs will be deposited on Zenodo with a citable DOI upon paper acceptance.

\bibliographystyle{elsarticle-num}
\bibliography{references}

\end{document}
